\section[A particular case of cohomological properness]{A particular case of cohomological properness: divisors
with relative normal crossings}
\label{XIII.2}

\setcounter{subsection}{-1}

\subsection{}
\label{XIII.2.0}
Let $R$ be a discrete valuation ring with field of fractions
$K$ and $L$ an étale $K$-algebra; $L$ is then a direct product
of a finite number of fields $L_i$, where $L_i$ is an étale
extension of~$K$. If $L'_i$ denotes the Galois extension
generated by $L_i$ in an algebraic closure of~$L_i$,
one says that $L$ is \emph{tamely ramified}
over~$R$ if the $L'_i$ are tamely
ramified extensions in the sense of~X~\Ref{X.3}, \ie if an inertia group~$I_i$ of~$L'_i|K$ is of order prime to the residue
characteristic $p$ of~$R$.

We know that $I_i$ is in any case an extension of a cyclic group
of order prime to $p$ by a $p$-group. (This follows from \cite[ch\ptbl IV prop\ptbl 7 cor\ptbl 4]{XIII.5} when the residue extension of~$R$ is
separable. The proof given in
\loccit extends to the general case in the following
way. Let us take up again the hypotheses and the notation of
\loccit but without assuming the residue extension
separable. Let $H_i$ be the subgroup of the inertia group $G_0$,
the set of the elements $s$ of~$G_0$ such that one has $s \pi /
\pi \in U^i$ for every uniformizer $\pi$ of~$A_L$. One then checks
that $G_0/H_1$ is a group of order prime to $p$ and that,
for $i \geqslant 1$, the $H^i/H^{i+1}$ are $p$-groups, from which
one deduces the announced result.)

\subsubsection{}
\label{XIII.2.0.1}
In the case where $R$ is strictly local, one has the following
simple characterization: the $K$-algebra $L$ is
tamely ramified over~$R$ if and only if the
$[L_i:K]$ are prime to $p$. Moreover, if $L$ is
tamely ramified over~$R$, the $L_i$ are cyclic extensions
of~$K$. Indeed, when $R$ is strictly local, $I_i$
is equal to the Galois group
% original p. 370
of~$L'_i$ over~$K$. As we have just recalled, $I_i$ is an extension
of a cyclic group of order prime to $p$ by a $p$-group. If
one assumes $L'_i$ tamely ramified over~$R$, $I_i$ is
then a cyclic group of order prime to $p$. It follows
that $[L_i:K]$ is prime to $p$ and that one has $L_i =
L'_i$. Conversely, if $[L_i:K]$ is prime to $p$, $I_i$ cannot
contain a nontrivial normal $p$-subgroup; $I_i$ is therefore
a cyclic group of order prime to $p$, which proves that $L_i$
is tamely ramified over~$R$.

\subsubsection{}
\label{XIII.2.0.2}
Let $R$ be a discrete valuation ring with field of fractions
$K$, $L$ an étale $K$-algebra, and let $\overline{R}$ be a
strict localization of~$R$, $\overline{K}$ its field of fractions,
$\overline{L} = L \otimes_K \overline{K}$. Then, for $L$ to be
tamely ramified over~$R$, it is necessary and sufficient that
$\overline{L}$ be tamely ramified over~$\overline{R}$. Indeed one reduces to the case where $L$ is a
field. Let then $\overline{L} = \prod_i \overline{L}_i$, where
the $\overline{L}_i$ are fields, extensions of~$\overline{K}$; if
$L'$ is the Galois extension generated by $L$, and, if
$\overline{L'} = L' \otimes_K \overline{K}$,
one likewise has a decomposition of~$\overline{L'}$ into a product
of fields,
$\overline{L'} = \prod_j \overline{L'_j}$
and each $\overline{L}_i$ is a subextension of at least one of the
$\overline{L'_j}$. Since $L'$ is a Galois extension of~$K$,
the $\overline{L'_j}$ are Galois extensions of
$\overline{K}$. Suppose $L$ tamely ramified over~$R$;
since the Galois group of~$\overline{L'_j}|K$ is isomorphic to the
inertia group of~$L'|K$, the $L'_j$ are also tamely
ramified over~$\overline{R}$, and hence so are the
$\overline{L}_i$. Conversely, suppose $\overline{L}$
tamely ramified over~$\overline{R}$. For each $j$,
let $v_j$ be the discrete valuation of~$\overline{L'_j}$ which extends
the valuation of~$\overline{K}$, and let us still denote by $v_j$ the valuation
induced on~$L'$. As $j$ varies, $v_j$ runs through the set of
valuations of~$L$ which extend the valuation of~$K$. Let $G =
\Gal(L'|K)$, $H=\Gal(L'|L)$, $I_j$ the inertia group of~$L'|K$ at
$v_j$, $J_j$ the inertia group of~$L'|L$ at $v_j$. The group $I_i$
is an extension of a cyclic group of order prime to $p$ by a
$p$-group $P_j$. Since the $\overline{L}_i$ are tamely
ramified over~$\overline{R}$, $I_j/J_j$ is of order prime to
$p$, hence one has $P_j \subset J_j$. Consequently the group $H$ contains
all the $P_j$, hence also the group $P$ generated by
% original p. 371
the $P_j$ for variable $j$. But the group~$P$ is invariant in $G$,
since an inner automorphism of~$G$ transforms the~$I_j$ into one
another, hence also the $P_j$ into one another. It follows that $P$ is a
subgroup of~$H$ normal in~$G$, hence, since $L'$ is
the Galois extension generated by $L$, that one has $P=1$, which
proves that $L$ is tamely ramified over~$R$.

More generally, let $R \to R'$ be a morphism of discrete
valuation rings such that the image of a uniformizer $\pi$ of
$R$ is a uniformizer $\pi'$ of~$R'$ and such that the residue
extension $\kres(R')$ is a separable extension of
$\kres(R)$. Let $K$ be the field of fractions of~$R$, $K'$ the field of
fractions of~$R'$, $L$ an étale $K$-algebra, $L' = L
\otimes_K K'$. Then, for $L$ to be tamely
ramified over~$R$, it is necessary and sufficient that $L'$ be
tamely ramified over~$R'$. Indeed we may assume
$R$ and $R'$ strictly local. By~\Ref{XIII.2.0.1} it suffices
to prove that, when $L$ is a field, so is~$L'$. Let $\tilde{R}$ be the normalization of~$R$ in $L$,
$\tilde{\pi}$ a uniformizer of~$\tilde{R}$, $\tilde{R}' =
\tilde{R} \otimes_R R'$. The extension $\kres(\tilde{R})|\kres(R)$ being
radicial and the extension $\kres(R')|\kres(R)$ étale,
$\kres(\tilde{R})\otimes_{\kres(R)}\kres(R')$ is a field [EGA IV 4.3.2 and
4.3.5]. This proves that $R'$ is a local ring, and, since $\pi$ has
image $\pi'$ in $R'$, one has $\kres(\tilde{R}')= R'/(\tilde{\pi})$;
consequently $R'$ is a discrete valuation ring \cite[ch\ptbl I \S 2
prop\ptbl 2]{XIII.5}, hence $L'$ is a field.

\subsubsection{}
\label{XIII.2.0.3}
By reduction to the strictly local case, one sees that a
subalgebra of a tamely ramified algebra is
tamely ramified, that the tensor product of two
tamely ramified algebras is tamely
ramified, that a tamely ramified algebra remains
so after extension of the discrete valuation ring,
that an algebra which becomes tamely ramified
after a tamely ramified extension is
tamely ramified.

\subsection{}
\label{XIII.2.1}
Let
% original p. 372
$X$ be an $S$-scheme, $D$ a divisor $\geq 0$ on~$X$.
Recall (SGA~5 II~4.2) that one says that $D$ has strictly
\emph{normal crossings relative to}~$S$
if there exists a finite family $(f_i)_{i\in I}$ of elements of
$\Gamma(X,\cal{O}_X)$ such that one has $D=\sum_{i\in
I}\divisor(f_i)$ and such that the following condition is satisfied:

\setcounter{subsubsection}{-1}
\subsubsection{}
\label{XIII.2.1.0}
For every point $x$ of~$\Supp D$, $X$ is smooth over~$S$ at $x$, and, if
one denotes by $I(x)$ the set of the $i\in I$ such that $f_i(x)=0$, the
subscheme $V((f_i)_{i\in I(x)})$ is smooth over~$S$ of
codimension
$\card I(x)$
in~$X$.

The divisor $D$ is said to have \emph{normal crossings relative
to} $S$ if, locally on~$X$ for the étale topology, it has
strictly normal crossings.

Let $D$ be a divisor with normal crossings relative
to~$S$. We set $Y=\Supp D$, $U=X-Y$, and we denote by $i\colon U\to X$
the canonical immersion. For every geometric point
$\overline{s}$ of~$S$ and for every maximal point $y$ of the geometric
fiber $Y_{\overline{s}}$, the ring
$R=\cal{O}_{X_{\overline{s}},y}$ is a discrete valuation
ring.

In the rest of this no., we shall use the following
technical definition:

\begin{subdefinition}
\label{XIII.2.1.1}
Let $F$ be a sheaf of sets on~$U$. One says that $F$ is
\emph{tamely ramified over}~$X$
(\emph{along}~$D$) \emph{relative to}~$S$ if, for every
geometric point $\overline{s}$ of~$S$, the following condition
is satisfied:

For every maximal point $y$ of~$Y_{\overline{s}}$, the restriction of
$F$ to the field of fractions $K$ of~$\cal{O}_{X_{\overline{s}}}$ is
representable by the spectrum of an étale $K$-algebra $L$,
tamely ramified over~$\cal{O}_{X_{\overline{s},y}}$.
\end{subdefinition}

Most often, when no confusion can result from it, we
shall omit the mention of~$D$ in the terminology.

\begin{subdefinition}
\label{XIII.2.1.2}
If
% original p. 373
$F$ is a sheaf of groups on~$U$, tamely
ramified over~$X$ relative to~$S$, one denotes by
$$
\H^1_{\tame}(U,F)
$$
the subset of~$\H^1(U,F)$ formed by the classes of torsors under
$F$ which are tamely ramified over~$X$ relative
to~$S$.
\end{subdefinition}

Let
$$
\xymatrix{ U \ar[r]^i \ar[d]_g & X \ar[dl]^f \\ T}
$$
be a commutative diagram of~$S$-schemes, with $i$ as
in~\Ref{XIII.2.1}; one denotes by
$$
\R^1_{\tame} g_* F
$$
the sheaf on~$T$ associated with the presheaf
$T'\mto\H^1_{\tame}(U',F)$, where $T'$ runs through the schemes
étale over~$T$ and where $U'=U\times_T T'$; $\R^1_{\tame} g_*F$
is a subsheaf of~$\R^1g_*F$.

Note that, if $g$ is coherent, if $\overline{t}$ is a geometric
point of~$T$, $\overline{T}$ the strict localization of
$T$ at $\overline{t}$, $\overline{U}=U\times_T\overline{T}$, one has an
isomorphism
\begin{equation*}
\label{eq:XIII.2.1.2.1}
\tag{\thesubsubsection.1} (\R^1_{\tame} g_*F)_{\overline{t}}\simeq
\H^1_{\tame}(\overline{U},\overline{F})
\end{equation*}

\subsubsection{}
\label{XIII.2.1.3}
Let $C_t((U,X)/S)$, or simply $C_t$,
be the category of étale coverings of~$U$ which are
tamely ramified over~$X$ relative
to~$S$. Suppose $U$ connected and let $a$ be a
geometric point of~$U$. Let $\Gamma_t$ be the functor which, to an
étale covering~$U'$ of~$U$ tamely ramified
over~$X$ relative to~$S$, associates the set of the geometric
points of~$U'$ over~$a$. From~\Ref{XIII.2.0}
it follows that the pair $(C_t,\Gamma_t)$ satisfies axioms
($G_1$) to ($G_6$) of V~\Ref{V.4}. Consequently $\Gamma_t$ is
representable by a pro-object which is called the
\emph{tamely ramified universal covering
% original p. 374
of~$(U,X)$ relative to~$S$ pointed at~$a$}.
The group opposite to the group of $U$-automorphisms of the
tamely ramified universal covering is called
the \emph{tamely ramified fundamental group}
and denoted
$$
\pi_1^{\tame}((U,X)/S,a)\text{ or simply }\pi_1^{\tame}(U,a)\text{ or
even }\pi_1^{\tame}(U).
$$
It is obviously a quotient of the fundamental group $\pi_1(U,a)$
(V~\Ref{V.6.9}).

\subsubsection{}
\label{XIII.2.1.4}
Let $F$ be a sheaf of groups on~$U$, $P$ a right torsor
with group $F$, $Q$ a left torsor with group $F$, and suppose
$P$ and $Q$ tamely ramified over~$X$ relative to~$S$. Then so is $P\stackrel{F}{\wedge}Q$. Indeed one
reduces to showing that, if $R$ is a discrete
valuation ring with field of fractions $K$ and if~$F$ is a
finite étale group scheme over~$K$, and if $P$ and $Q$ are
two torsors under~$F$ tamely ramified over~$R$, then
so is $P\stackrel{F}{\wedge} Q$. Now
$T=P\stackrel{F}{\wedge} Q$ is a quotient of~$P\times_K Q$. If $L$,
$M$, $N$ denote the $K$-algebras representing
$T$, $P$, $Q$ respectively, then $L$ is a subalgebra of
$M\otimes_K N$, and it follows from~\Ref{XIII.2.0.3} that $L$ is
tamely ramified over~$R$.

We deduce from the preceding that, if $F$ is a sheaf of
groups on~$U$ and if there exists a torsor with group $F$
tamely ramified over~$X$ relative to~$S$, then
$F$ is tamely ramified over~$X$ relative to~$S$. Indeed the torsor $P^{\circ}$ opposite to~$P$ is
tamely ramified over~$X$ relative to~$S$,
since it is isomorphic to $P$ as a sheaf of sets. If
$\leftexp{P\mkern-4mu}{F}$ is the group obtained by twisting $F$ by $P$, one has an
isomorphism
$$
F\simeq P^{\circ}\stackrel{\leftexp{P\mkern-4mu}{F}}{\wedge}P
$$
and consequently $F$ is tamely ramified over~$X$
relative to~$S$.

One sees as before that, if $F\to F'$ is a morphism
of sheaves of groups on~$U$, tamely ramified over~$X$ relative to~$S$, and if $P$ is a torsor under $F$
tamely ramified over~$X$ relative
% original p. 375
to~$S$, then the torsor $P'$ deduced from~$P$ by the extension of the
structure group $F\to F'$ is tamely ramified over~$X$
relative to~$S$.

In particular the canonical morphism
$$
\H^1(U,F)\to \H^1(U,F')
$$
gives by restriction to $\H^1_{\tame}(U,F)$ a canonical morphism
$$
\H^1_{\tame}(U,F)\to\H^1_{\tame}(U,F')
$$

\subsubsection{}
\label{XIII.2.1.5}
Let $S'\to S$ be a morphism and let us denote by $U'$ (\resp $X'$, etc.)
the inverse image of~$U$ (\resp $X$, etc.) over~$S'$. If $F$ is a
sheaf of sets on~$U$ tamely ramified over~$X$
relative to~$S$, it follows from
definition~\Ref{XIII.2.1.1} and from~\Ref{XIII.2.0.3} that $F'$ is
tamely ramified over~$X'$ relative to~$S'$.

If now $F$ is a sheaf of groups on~$U$, the inverse image
over~$S'$ of a torsor under $F$ tamely ramified over~$X$
relative to~$S$ is a torsor under $F'$ tamely
ramified over~$X'$ relative to~$S'$. In particular one has a
canonical functor
\begin{equation*}
\label{eq:XIII.2.1.5.1}
\tag{\thesubsubsection.1} C_t((U,X)/S)\to C_t((U',X')/S').
\end{equation*}
Suppose $U$ and $U'$ connected and let $a$ be a
geometric point of~$U$, $a'$ a geometric point of~$U'$
over~$a$; one deduces from the preceding a canonical
morphism
\begin{equation*}
\label{eq:XIII.2.1.5.2}
\tag{\thesubsubsection.2} \pi_1^{\tame}(U',a')\to\pi_1^{\tame}(U,a).
\end{equation*}

If $S'\to S$ is a morphism and $h\colon T'\to T$ the canonical
projection, the morphism
$$
h^*(\R^1g_*F)\to\R^1g_*'F'
$$
gives by restriction a canonical morphism
\begin{equation*}
\label{eq:XIII.2.1.5.3}
\tag{\thesubsubsection.3} h^*(\R^1_{\tame}g_*F)\to\R^1_{\tame}g_*'F'
\end{equation*}

\subsubsection{}
\label{XIII.2.1.6}
Let
% original p. 376
$F$ be a sheaf of groups on~$U$, tamely
ramified over~$X$ relative to~$S$. The notation being
that of~\Ref{XIII.2.1.2}, one has canonical exact sequences:
\begin{equation*}
\label{eq:XIII.2.1.6.1}
\tag{\thesubsubsection.1}
\begin{array}{ccccccc}
1 & \longrightarrow & \H^1(X,i_*F)& \longrightarrow & \H^1_{\tame}(U,F) &
\longrightarrow & \H^0(X,\R^1_{\tame} i_*F)\\[10pt] 1 &
\longrightarrow & \R^1f_*(i_*F) & \longrightarrow & \R^1_{\tame}g_*F &
\longrightarrow & f_*(\R^1_{\tame}i_*F).\end{array}
\end{equation*}

The first is obtained from the exact sequence
(SGA~4 III~3.2):
$$
1\to \H^1(X,i_*F)\to \H^1(U,F)\to
\H^0(X,\R^1i_*F).
$$
Indeed it suffices to show that the image of~$\H^1(X,i_*F)$ in
$\H^1(U,F)$ is in fact contained in $\H^1_{\tame}(U,F)$ and that the image of
$\H^1_{\tame}(U,F)$ in $\H^0(X,\R^1i_*F)$ is in fact contained in
$\H^0(X,\R^1_{\tame}i_*F)$. Now the inverse image over~$U$ of a torsor under
$i_*F$ is a torsor under $i^*i_*F$ which is obviously
tamely ramified over~$X$ relative to~$S$, hence the same
holds after the extension of the structure group
$i^*i_*F\to F$, which proves the existence of the arrow
$\H^1(X,i_*F)\to\H^1_{\tame}(U,F)$. The fact that the image of~$\H^1_{\tame}(U,F)$
in $\H^0(X,\R^1i_*F)$ is contained in $\H^0(X,\R^1_{\tame}i_*F)$
follows immediately from the definition of~$\R^1_{\tame}i_*F$. This
proves the existence of the first exact sequence, and the second
is deduced from it by localization.

\subsection{}
\label{XIII.2.2}
We keep the notation of~\Ref{XIII.2.1}. We are going to define
a notion of tamely ramified object of a stack $\Phi$
on~$U$ when the latter is given, locally \emph{on} $X$ and
on~$S$ for the étale topology, as \emph{the inverse image of a
stack} $\Psi$ \emph{on}~$S$.

First let $G$ be a gerbe on~$U$ and suppose given an
étale surjective morphism $S_1\to S$, an étale surjective morphism
$X_2\to X\times_S S_1$, a trivial gerbe $H$ on~$S_1$ and an
isomorphism
$$
G|{U_2}\to H|{U_2},
$$
where
% original p. 377
$U_1=U\times_X X_1$, $U_2=U\times_X X_2$. When one chooses a
trivialization of~$H|{X_2}$, the isomorphism above identifies
$G|{U_2}$ with the stack of torsors under a sheaf of groups $F$. One
says that an element $x$ of~$G_U$ is tamely
ramified over~$X$ relative
to~$S$
if the restriction of~$x$
to $U_2$ is a torsor tamely ramified over~$X$
relative to~$S$. By~\Ref{XIII.2.1.4} this notion does not
depend on the way in which one has trivialized~$H|X_2$.

Now let $\Phi$ be a stack on~$U$ and suppose given an
étale surjective morphism $S_1\to S$, an étale surjective
morphism $X_2\to X\times_S S_1$, a stack $\Psi$ on~$S_1$ and an
isomorphism
$$
i\colon \Phi|U_2\to \Psi|U_2.
$$
Let $x$ be an element of~$\Phi_U$, $G_x$ the maximal subgerbe
of~$\Phi$ generated by $x$
\cite[III~2.1.7]{XIII.1}, $S\Phi$ the sheaf of
maximal subgerbes of~$\Phi$. The isomorphism $i$ induces an
isomorphism
$$
S\Phi|U_2\to S\Psi|U_2.
$$
It follows from~\Ref{XIII.5.7} that, after replacing $S_1$ if necessary by
an étale surjective extension, one has a unique maximal
subgerbe $H$ of~$\Psi$, which one may assume trivial, such that $i$
defines an isomorphism
$$
G_x|U_2\to H|U_2.
$$
One says that the element $x$ is tamely ramified
over~$X$ relative to~$S$ if it is so as an element
of~$G_x$ endowed with the isomorphism above.

\subsubsection{}
\label{XIII.2.2.1}
Let $\Phi$ be a stack on~$U$ given, locally on~$X$ and $S$,
as the inverse image of a stack on~$S$, and let
$$
\xymatrix{ U \ar[r]^i
\ar[d]_g & X \ar[dl]^f \\ T }
$$
be a diagram as in~\Ref{XIII.2.1.2}. For every scheme $T'$
étale over~$T$, if $U'=U\times_T T'$, one considers the
subset $(g_*^{\tame}\Phi)_{T'}$
of~$(g_*\Phi)_{T'}=\Phi_{U'}$
formed
% original p. 378
by the elements of~$\Phi_{U'}$ which are tamely
ramified
over~$X$ relative to~$S$. One calls tamely ramified
direct image
of~$\Phi$ by~$g$, and one denotes by
$$
g_*^{\tame}\Phi
$$
the full subcategory of~$g_*\Phi$ whose objects over
a scheme $T'$ étale over~$T$ are the elements of
$(g_*^{\tame}\Phi)_{T'}$. It is clear that $g_*^{\tame}\Phi$ is a substack of
$g_*\Phi$.

\subsubsection{}
\label{XIII.2.2.2}
By reduction to the case of a stack of torsors, one sees that, if
$\Phi$ is a stack on~$U$ which is, locally for the
étale topology of~$S$ and $X$, the inverse image of a stack on~$S$, the canonical
morphism
$$
h^*(g_*\Phi)\to g_*'\Phi'
$$
gives by restriction a canonical morphism
$$
h^*(g_*^{\tame}\Phi)\to g_*^{\prime\tame}\Phi'
$$

\begin{remarks}
\label{XIII.2.3}

a) If $F$ is a locally constant constructible sheaf of sets
on~$U$, for $F$ to be tamely
ramified over~$X$ relative to~$S$, it suffices that the condition
of~\Ref{XIII.2.1.1} be satisfied for the geometric
points of~$S$ over the maximal points of~$S$. To
see this we may assume that the divisor $D$ has strictly normal
crossings. The sheaf $F$ is representable by an étale
covering $V$ of~$U$. If $\overline{s}$ is a geometric
point of~$S$, $y$ a maximal point of
$Y_{\overline{s}}$, we denote by $\overline{S}$ the strict localization of
$S$ at $\overline{s}$, $\overline{X}$ the strict localization of~$X$ at
$\overline{y}$, $\overline{U}=U\times_X \overline{X}$,
$\overline{V}=V\times_X\overline{X}$. If the condition
of~\Ref{XIII.2.1.1} is satisfied at the geometric points
over the maximal points of~$S$, it follows from~\Ref{XIII.5.5}
below that $\overline{V}$ is a covering of~$\overline{U}$
tamely ramified over~$\overline{X}$ relative to
$\overline{S}$; consequently $V$ is an étale covering of~$U$
tamely ramified over~$X$ relative to~$S$.

b)
% original p. 379
Let $F$ be a sheaf of groups on~$U$ tamely
ramified over~$X$ relative to~$S$. If $\overline{s}$ is a
geometric point of~$S$, $y$ a maximal point of
$Y_{\overline{s}}$, we denote by $K$ the field of fractions of
$\cal{O}_{X_{\overline{s}},y}$. Suppose that, for every point
$\overline{s}$ and for every point $y$, the $K$-algebra $L$ whose
spectrum represents $F|K$ is of rank prime to the
residue characteristic $p$ of
$\cal{O}_{X_{\overline{s}}}$. We shall sometimes say, by abuse of language,
that $F$ is prime to the residue characteristics of~$S$.
When this is so, every torsor $P$ under $F$ is
tamely ramified over~$X$ relative to~$S$.

Indeed, let $\overline{R}$ be the strict localization of
$\cal{O}_{X_{\overline{s}},y}$ at $\overline{y}$, $\overline{K}$ its
field of fractions, $\overline{F}$ the inverse image of~$F$ on~$\overline{K}$. Let us show that we may assume $F$ constant. Since
$F$ is tamely ramified over~$X$ relative to~$S$,
$\overline{F}$ is representable by the spectrum of a
$\overline{K}$-algebra $L=\prod L_i$, where the $L_i$ are
extensions of~$\overline{K}$ of degree prime to $p$. One can
therefore find an extension $K'$ of~$\overline{K}$ of degree prime
to $p$ such that $\overline{F}|K'$ is a constant
sheaf. By~\Ref{XIII.2.0.3}, to prove that
$P|\overline{K}$ is tamely ramified over~$\overline{R}$, it suffices to see that $P|K'$ is tamely
ramified over the integral closure of~$\overline{R}$ in
$K'$, whence the reduction to the case where $\overline{F}$ is
constant. Suppose henceforth $\overline{F}$ constant. The
$\overline{K}$-algebra $H$ which represents $P|\overline{K}$ is
then a product of extensions $H_i$ of~$\overline{K}$ isomorphic to
one another. Since the rank of~$H$ is prime to $p$, so is
$[H_1:\overline{K}]$, which proves that $H$ is
tamely ramified over~$X$ relative to~$S$.

c) Let $X$ be a regular scheme, $D$ a divisor with normal
crossings of~$X$ (SGA~5 I~3.1.5), $U=X-\Supp D$, $F$ a
sheaf of sets on~$U$. If $y$ is a maximal point of~$\Supp
D$, we denote by $K$ the field of fractions of
$\cal{O}_{X,y}$. One says that $F$ is \emph{tamely
ramified relative to}~$D$
if, for every maximal point $y$ of~$\Supp D$, $F|K$ is
representable by a $K$-algebra tamely
ramified over~$\cal{O}_{X,y}$.
\end{remarks}
